% =====================================================================
% 陈炳林回忆录   印刷版
% =====================================================================
% 编译方法：  xelatex book.tex   (跑两遍以生成目录)
% 字体：     Noto Serif CJK SC (思源宋体) + Noto Sans CJK SC (思源黑体)
% 版式：     A5  140mm × 210mm   ctexbook  传统中文书籍版式
% =====================================================================

\documentclass[a5paper,11pt,openright,oneside]{ctexbook}

% ---------- 页面 ----------
\usepackage[
  paperwidth=148mm, paperheight=210mm,
  inner=20mm, outer=18mm, top=22mm, bottom=24mm,
  headsep=8mm, footskip=14mm,
  heightrounded
]{geometry}

% ---------- 字体 ----------
\usepackage{fontspec}
\usepackage{xeCJK}

% 西文字体：Latin Modern Roman + Latin Modern Sans
\setmainfont{Latin Modern Roman}
\setsansfont{Latin Modern Sans}

% 中文字体：思源宋体（正文）/ 思源黑体（标题）
% 直接用文件路径加载（NotoSerif/Sans CJK 是 .ttc 集合）
\setCJKmainfont[
  Path           = /usr/share/fonts/opentype/noto/,
  Extension      = .ttc,
  UprightFont    = NotoSerifCJK-Regular,
  BoldFont       = NotoSerifCJK-Bold,
  ItalicFont     = NotoSerifCJK-Light,
  BoldItalicFont = NotoSerifCJK-SemiBold
]{NotoSerifCJK-Regular}

\setCJKsansfont[
  Path           = /usr/share/fonts/opentype/noto/,
  Extension      = .ttc,
  UprightFont    = NotoSansCJK-Light,
  BoldFont       = NotoSansCJK-Bold
]{NotoSansCJK-Light}

% 备用：思源宋体 Light 当楷体用（ctex 已有 \kaishu/\booksans，所以换名字避免冲突）
\newCJKfontfamily\bookkai[
  Path           = /usr/share/fonts/opentype/noto/,
  Extension      = .ttc,
  UprightFont    = NotoSerifCJK-Light
]{NotoSerifCJK-Light}

\newCJKfontfamily\booksans[
  Path           = /usr/share/fonts/opentype/noto/,
  Extension      = .ttc,
  UprightFont    = NotoSansCJK-Light,
  BoldFont       = NotoSansCJK-Bold
]{NotoSansCJK-Light}

\newCJKfontfamily\booksansthin[
  Path           = /usr/share/fonts/opentype/noto/,
  Extension      = .ttc,
  UprightFont    = NotoSansCJK-Thin
]{NotoSansCJK-Thin}

% 数字与英文字号：用衬线 sans 风格的 Latin Modern Sans Thin（无 -> 用普通拉丁数字）
\newfontfamily\figurefont{Latin Modern Sans}

% ---------- 行距与段落 ----------
\linespread{1.55}                 % 中文书行距 1.5-1.6 倍较舒服
\setlength{\parindent}{2em}       % 段首缩进两字
\setlength{\parskip}{2pt plus 1pt}

% ---------- 章节标题样式（核心：高对比 typography）----------
\usepackage{titlesec}
\usepackage{xcolor}
\definecolor{inkblack}{HTML}{1a1a1a}
\definecolor{inkgray}{HTML}{595959}
\definecolor{inkfaint}{HTML}{9c9c9c}
\definecolor{seal}{HTML}{8B2C1A}     % 印章红，留作装饰

% 章首页：极简风。上方一行极细 "CHAPTER" 数字 sans，
% 中部超大号 sans-thin 数字，下方一条短线，再下方大号宋体章名，
% 再下方副题（年代）和长短线。size jump ~3x.
\titleformat{\chapter}[display]
  {\thispagestyle{plain}\vspace*{1.6cm}\centering}
  {%
    {\booksansthin\fontsize{12pt}{14pt}\selectfont\color{inkfaint}\addfontfeatures{LetterSpace=20} \MakeUppercase{\chaptername}}\\[6pt]
    {\booksansthin\fontsize{84pt}{84pt}\selectfont\color{inkblack}\thechapter}\\[18pt]
    {\color{inkfaint}\rule{2em}{0.4pt}}\\[20pt]
  }
  {0pt}
  {\fontsize{24pt}{30pt}\selectfont\bfseries\color{inkblack}}
  [\vspace{12pt}{\color{inkfaint}\rule{2em}{0.4pt}}\vspace{2.4cm}]

% 节标题
\titleformat{\section}
  {\booksans\Large\bfseries\color{inkblack}}
  {\thesection}{0.8em}{}

\titleformat{\subsection}
  {\booksans\large\color{inkblack}}
  {}{0em}{}

% ---------- 页眉页脚 ----------
\usepackage{fancyhdr}
\pagestyle{fancy}
\fancyhf{}
\renewcommand{\headrulewidth}{0pt}
% 页眉：左中右皆空；中部放章名，前后各一个小菱点修饰
\fancyhead[C]{%
  \booksansthin\footnotesize\color{inkfaint}%
  \raisebox{2pt}{\tiny ◇}\hspace{0.8em}%
  \color{inkgray}\leftmark%
  \hspace{0.8em}\raisebox{2pt}{\color{inkfaint}\tiny ◇}%
}
\fancyfoot[C]{\booksansthin\small\color{inkgray}\thepage}

\fancypagestyle{plain}{%
  \fancyhf{}
  \fancyfoot[C]{\booksansthin\small\color{inkgray}\thepage}
  \renewcommand{\headrulewidth}{0pt}
}

\renewcommand{\chaptermark}[1]{\markboth{#1}{}}

% ---------- 目录样式 ----------
\usepackage{titletoc}
\usepackage{tocloft}
\renewcommand{\cftchapfont}{\normalfont}
\renewcommand{\cftchappagefont}{\booksansthin\color{inkgray}}
\renewcommand{\cftchapleader}{\booksansthin\color{inkfaint}\cftdotfill{\cftdotsep}}
\renewcommand{\cftdotsep}{2.0}
\setlength{\cftbeforechapskip}{10pt}
\setlength{\cftchapindent}{0pt}
\setlength{\cftchapnumwidth}{2.6em}

% 关掉 ctexbook 默认的 "第 X 章" 章节编号显示，由我们自定义
% （已在 titleformat 里处理）

% ---------- 图片 ----------
\usepackage{graphicx}
\usepackage[export]{adjustbox}
\usepackage{enumitem}
\graphicspath{{images/}}

% 安全图片：文件存在则插入，否则用占位框
\newcommand{\safeimage}[1]{%
  \IfFileExists{images/#1}{%
    \includegraphics[max width=0.85\linewidth,max totalheight=0.5\textheight,keepaspectratio]{#1}%
  }{%
    \fbox{\parbox{0.75\linewidth}{\centering\booksansthin\small\color{inkfaint}%
      \vspace{1.5em}[\,待补图：\detokenize{#1}\,]\vspace{1.5em}}}%
  }%
}

% ---------- 引文样式 ----------
\usepackage{quoting}
\quotingsetup{
  leftmargin=2em,
  rightmargin=2em,
  vskip=0.4em
}
\renewenvironment{quote}
  {\begin{quoting}\bookkai\small}
  {\end{quoting}}

% ---------- 链接与目录支持 ----------
\usepackage[hidelinks,unicode,bookmarks=true,bookmarksnumbered=true]{hyperref}
\hypersetup{
  pdftitle={陈炳林回忆录},
  pdfauthor={陈炳林  口述／陈旭鹏  整理},
  pdfsubject={家族回忆录},
  pdfkeywords={陈炳林,回忆录,家史,memoir}
}

% ---------- 一些排版小工具 ----------
% \usepackage{soul}      % 不需要
% \usepackage{microtype} % xelatex 下功能有限，省略

% 章前题记（章名下方的副标题或时间提示）
\newcommand{\chapterepigraph}[1]{%
  \begin{center}\booksansthin\small\color{inkgray}#1\end{center}\vspace{1em}}

% 分节装饰符
\newcommand{\dinkus}{\par\vspace{1em}\begin{center}\color{inkfaint}❦\end{center}\vspace{1em}\par}

% horizontal rule -> 用 dinkus 替代
\let\oldhrule\hrule

% ---------- 关闭 ctexbook 默认章节编号样式 ----------
% ctexbook 默认会输出 "第一章 章名"，我们的 titleformat 已经独立处理章号
\CTEXsetup[number={\arabic{chapter}},name={,}]{chapter}

% ---------- 书末杂项 ----------
\usepackage[absolute,overlay]{textpos}

% =====================================================================
\begin{document}
% =====================================================================

% --------- 半扉页 / Half-title ---------
% 仅有书名，极简，象征"翻开第一页"的仪式感
\pagestyle{empty}
\thispagestyle{empty}
\vspace*{6cm}
\begin{center}
  {\booksansthin\fontsize{18pt}{22pt}\selectfont\color{inkgray} 陈炳林回忆录}
\end{center}
\cleardoublepage

% --------- 扉页 / Title page ---------
\begin{titlepage}
  \centering
  \vspace*{2cm}
  {\booksansthin\fontsize{20pt}{24pt}\selectfont\color{inkfaint}\addfontfeatures{LetterSpace=30} MEMOIR}\\[6pt]
  {\booksansthin\fontsize{14pt}{18pt}\selectfont\color{inkfaint}回\quad 忆\quad 录}\\[2.6cm]
  {\fontsize{52pt}{60pt}\selectfont\bfseries 陈\,炳\,林}\\[0.8cm]
  {\color{inkfaint}\rule{4em}{0.4pt}}\\[0.6cm]
  {\booksansthin\large\color{inkgray}一九四〇\,—\,二〇二六}\\[2.8cm]
  {\booksansthin\small\color{inkgray}陈炳林\ 口述\quad 陈旭鹏\ 整理}
\end{titlepage}

\cleardoublepage
\thispagestyle{empty}
\vspace*{\fill}
\begin{center}
  {\bookkai\fontsize{16pt}{22pt}\selectfont 谨以此书\\[8pt]纪念我的爷爷}\\[2cm]
  {\color{inkfaint}\rule{2em}{0.3pt}}\\[1.6cm]
  {\booksansthin\small\color{inkgray}并献给\\[6pt]爱他的家人}
\end{center}
\vspace*{\fill}
\cleardoublepage

% --------- 前言：追思序 ---------
\frontmatter
\pagestyle{fancy}

\chapter*{追思序}
\addcontentsline{toc}{chapter}{追思序}
\markboth{追思序}{}

% 追思序草稿  —  请按需修改 / 替换
% 编辑这一段就行；正文紧接 \chapter*{追思序} 之后

\vspace{1em}

二〇二六年初夏，爷爷走完了他八十六年的人生。这本小册子，是他亲口讲述的回忆，是他亲笔写下的童年，是我们一家人为他做的一件晚到很多年的事。

爷爷讲故事的时候，总是从一九四〇年的春天开始。那是他出生的年份，也是河南灾年的前夕。他记得家里穷得只剩两间草房；记得母亲生病、父亲徒步几日换回的两升麸面救了他的命；记得祖母带着大伯远走他乡，留下他在乡邻黄奶的饭碗里慢慢长大。这些事他讲过很多遍，每次讲，神情都和第一次一样郑重。

二〇一七年我还在念书的时候，曾想替他把这些故事整理成一本小书。那时只整理出十几篇，权且作了一份课程作业，又随手放进我的博客里。没想到这成了它现存唯一完整的版本。原本以为日后还有大把时间，可以一篇篇补、一段段问、把家里所有人的名字、所有的年份、所有他想说而未说的，慢慢都收进来。

现在没有这样的时间了。

但是他写下的那些字还在。那些关于一九五八年的灾难、关于他和奶奶相识的初中校园、关于工作时第一次远行的清晨、关于在县委大院里熬夜写材料的笔尖——都还在。把它们重新认认真真地排好版，印成一本真正的书，让家里每个人都能拿到一册，让以后我们的孩子也能读到，是我现在唯一还能为他做的事情。

这本书保留了他原本的口吻。文字几乎没有改动——他的语气、他的用词、他偶尔的跳跃和重复，都还是他的。我们能听到他坐在沙发上、戴着老花镜、就着茶杯讲故事的样子。这就够了。

爷爷，我把您的书整理好了。

\vspace{2em}
\begin{flushright}
{\bookkai 陈旭鹏\\[2pt]\small 二〇二六年五月}
\end{flushright}


\cleardoublepage

% --------- 原序 ---------
\chapter*{原序}
\addcontentsline{toc}{chapter}{原序}
\markboth{原序}{}

\input{src/00-preface-original-body}

\cleardoublepage

% --------- 目录 ---------
{\booksansthin\renewcommand{\contentsname}{目\hspace{1em}录}\tableofcontents}
\cleardoublepage

% --------- 正文：八章 ---------
\mainmatter

\input{src/02-chapter-1-body}
\input{src/03-chapter-2-body}
\input{src/04-chapter-3-body}
\input{src/05-chapter-4-body}
\input{src/06-chapter-5-body}
\input{src/07-chapter-6-body}
\input{src/08-chapter-7-body}
\input{src/09-chapter-8-body}

% --------- 附录 ---------
\appendix
\renewcommand{\thechapter}{\Roman{chapter}}

\chapter{童年的诗}

玩秋千 石人山情节

2016.9.11

茅草小屋金銮殿，桃杏梨花来装点。 春日争相露笑脸，葛花树上玩秋千。 若有同龄孩童伴，情景定会不一般。 良辰美景一人享，亦开心来亦孤单。

石人山景 石人山景美若画，儿时记忆难忘下。 高山顶峰衬其峻，云雾缭绕摹其神。 谷满林竹花锦绣，潺潺流水汇清泉。

恋家 少时寄住石人山，犹似修隐仙景间。 天亮登山陪桑蚕，入夜土炕头边玩。 山风怒吼夜狼嚎，柴房家犬牙齿露。 一日三餐尚果腹，思乡遥望东方愁。

四伯 追思堂伯父 幼时父母皆亡殁，凄凉处境人怜惜。 尝尽人间冷与暖，变得木讷而寡言。 不养鸡鸭不养猪，不种蔬菜只种谷。 厨中并无盐油醋，舌尖哪有美味留。

无题 民国三十年，恰逢吾降生。 上天不下雨，庄家未收成。 家园遭天灾，倭寇来侵略。 人在襁褓中，母亲患大病。 断了生命水，差点送性命。 - 注：生命之水指母亲奶水。

思往昔看今朝

2015.3

忆往昔峥嵘岁月稠， 想痛处，满脸沟壑老泪流。 看今朝国盛家富， 家和睦儿孙老人事事顺溜。

忆石人涧茅草小屋“金銮殿”

2016.9.9

石人身高八丈三，一半身在云雾间。 石婆闲来无去处，身贴悬崖避风寒。 谷底花树高又俊，秀拔挺立直冲天。 房后园中绿竹林，粗若碗口真稀罕。 门前小河水不断，源自上游山泉涧。 春风吹佛小山村，桃杏梨花齐斗艳。 吾辈此时何处去，葛藤树上荡秋千。

咏春 石人山旧居 茅草小屋金銮殿，桃杏梨花来装点。 三亩翠竹作盆景，门前溪水流不断。 欲问清水何处来，遥指山谷泉水涧。 春色美景独享有，葛花树上荡秋千。

大竹园变堰潭 旧时竹园改了观，现已变成大堰潭。 鱼儿为觅漂浮时，头儿伸出乱眨眼。 若想尝下新鲜味，一网下去捞两篮。

\begin{itemize}[leftmargin=2.5em,itemsep=2pt]
  \item 注：原来的大竹园在58年时，竹子被砍伐，修成了大堰潭，原生态的景观也被破坏殆尽
\end{itemize}


\input{src/11-appendix-justice-body}
\input{src/12-appendix-essays-body}
\chapter{年表}

\textbf{1940} 阴历七月初九出生于石桥街西夹后布袋街外祖母家的巷子里。辗转生活于薛庄、沙山、石人沟、魏庄。

\textbf{1941} 阴历五月十七妻子王若平出生于蒲山镇。

\textbf{1948} 患天花

\textbf{1949} 新中国成立

\textbf{1952} 土改，在魏庄分得三间大瓦房和几亩地

\textbf{1956} 读初中

\textbf{1959} 秋季，就读于郑州农业机械化专科学校（一年后更名为河南农学院农机分院），现为河南工学院。

\textbf{1961} 全国大中院校放长假，停课回家，在生产队种地。

\textbf{1962} 农历七月十一与王若平结婚，住在魏庄

\textbf{1963} 阴历五月初一，大儿子陈华军出生于魏庄，送至沽沱抚养。八月下旬与妻子重回郑州上学，九月一日报道。

\textbf{1964} 八月中旬毕业分配在茶庵工作，妻子分配至安皋，两地分居。

\textbf{1966} 四月，调动至地方国营南阳县拖拉机总站机务科（农机局前 身）

\textbf{1967} 阴历五月二十九，二儿子陈永出生于南阳地区中心医院。阴 历六月十三，儿媳杜思桂出生

\textbf{1968} 秋天，妻子担任拖拉机修配厂车工。

\textbf{1969} 阴历四月初三，三儿子陈斌于南阳市医院出生。

\textbf{1971} 调动至农业局工作，在李华庄药厂担任政工组长。阴历十一月十七，小儿子陈涛于南阳地区中心医院出生。

\textbf{1972} 在煤厂街盖房子，由妻子主要负责。

\textbf{1974} 妻子调动至南阳县计委仓库保管员，一边工作一边照顾四个儿子。

\textbf{1976} 调动至县委办公室，工作几个月时间。

\textbf{1976} 阴历十月十九，儿媳宋爱平出生。

\textbf{1977} 调至组织部工作。

\textbf{1978} 阴历十月初一，儿媳刘延娇出生。

\textbf{1980} 阴历三月二十九，儿媳张娟出生。

\textbf{1980} 阴历七月十二父亲陈锡海逝世，葬于皇路店魏庄

\textbf{1983} 调动至县人事局右派改正办公室工作，主要负责右派改正。

\textbf{1984} 调动至南阳县农机局监理站工作。

\textbf{1985} 腊月初一母亲吕金兰逝世，葬于皇路店魏庄

\textbf{1989} 阴历八八年腊月十九陈华军与范惠杰结婚、腊月二十陈永与 杜思桂结婚。

\textbf{1989} 阴历九月初八孙女陈旭鲲出生于南阳地区中心医院。

\textbf{1996} 妻子退休，55岁。单位为县物资局金属公司。

\textbf{1995} 阴历八月二十五孙子陈旭鹏出生于南阳地区中心医院。

\textbf{1999} 阴历二月二十六陈斌与宋爱平结婚。

\textbf{2000} 因身体原因退休，60岁。

\textbf{2001} 阴历正月初七孙女陈泓妤出生于南阳市县医院。

\textbf{2002} 阴历十一月初九陈涛与刘延娇结婚。

\textbf{2003} 阴历六月初二孙女陈奕霏出生于南阳市县医院。

\textbf{2004} 阴历三月初九陈华军与张娟结婚。 阴历十一月初四孙女陈佳欣出生于南阳市县医院。

\textbf{2012} 阴历五月十二孙女陈映竹出生于南阳地区中心医院。

\textbf{2016} 阴历一五年腊月十三孙子陈旭宏出生于南阳市县医院。

\textbf{2017} 阴历三月十九，重建老家魏庄的住宅。

\textbf{2017} 阴历九月初三孙女陈旭鲲与刘一佳结婚。

\textbf{2017} 年末，老家新房建成。

\textbf{2018} 阴历六月十八重孙女刘沁孜出生。

\textbf{2025} 5月25日，孙子陈旭鹏与孟舒晨结婚。

\textbf{2026} 5月24日15时28分，在南阳市中心医院安然离世，享年八十七岁。


\clearpage
\begin{landscape}
\thispagestyle{plain}
\null\vspace*{-0.4cm}
\begin{center}
\begin{tikzpicture}[
  remember picture,
  every node/.style={
    draw=seal, rounded corners=2.5pt, line width=0.45pt,
    align=center, font=\trueka\small,
    inner sep=3.5pt, minimum width=1.55cm
  },
  anc/.style={
    fill=seal!6, font=\trueka\small\bfseries,
    minimum width=1.85cm, inner sep=4.5pt
  },
  spouse/.style={
    fill=seal!10, font=\trueka\small\bfseries,
    minimum width=1.85cm, inner sep=4.5pt
  },
  small/.style={font=\trueka\footnotesize, minimum width=1.35cm, inner sep=2.8pt},
  ln/.style={draw=seal!60, line width=0.45pt},
  marriage/.style={draw=seal!75, line width=0.7pt, double, double distance=1pt},
  x=1cm, y=1cm,
]

% ── Columns ────────────────────────────────────────────
\def\cA{0}      % 祖辈
\def\cB{3.2}    % 炳林一代
\def\cC{6.4}    % 子女
\def\cD{9.5}    % 孙辈
\def\cE{12.6}   % 重孙辈

% ── 标题（嵌入 tikz 内部，确保与图同页） ─────────────
\node[draw=none, font=\trueka\fontsize{18pt}{22pt}\selectfont\bfseries, text=inkblack]
  (title) at (\cC, 7.2) {家\hspace{0.8em}族\hspace{0.8em}世\hspace{0.8em}系};
\node[draw=none, below=2pt of title]
  {{\color{seal}\rule{1.2em}{0.4pt}}\hspace{0.4em}%
    \raisebox{-1pt}{\color{seal}\small$\Diamond$}\hspace{0.4em}%
    {\color{seal}\rule{1.2em}{0.4pt}}};

% ── 第一代：双方父母 ──────────────────────────────────
\node[anc] (chenP) at (\cA, 1.6)  {陈\,锡\,海\\吕\,金\,兰};
\node[anc] (wangP) at (\cA, -2.3) {王\,铭\,鼎\\庆\hspace{1em}良};

% ── 第二代：炳林兄妹 + 若平 ──────────────────────────
\node (cby)  at (\cB,  4.2) {陈炳郁};
\node (cbq)  at (\cB,  3.3) {陈炳琴};
\node (cbi)  at (\cB,  2.4) {陈炳义};
\node (cbll) at (\cB,  1.5) {陈炳立};
\node (cbyu) at (\cB,  0.6) {陈炳玉};
\node[spouse] (cbl) at (\cB, -0.65) {陈\,炳\,林};
\node[spouse] (wrp) at (\cB, -2.30) {王\,若\,平};

% 婚姻线（双线表示婚配）
\draw[marriage] (cbl.south) -- (wrp.north);
\coordinate (union) at ($(cbl.south)!0.5!(wrp.north)$);

% 陈家父母 → 五位兄妹 + 炳林
\foreach \n in {cby,cbq,cbi,cbll,cbyu,cbl} {
  \draw[ln] (chenP.east) to[out=0,in=180] (\n.west);
}
% 王家父母 → 若平
\draw[ln] (wangP.east) to[out=0,in=180] (wrp.west);

% ── 第三代：炳林若平之子女 ────────────────────────────
\node (chj) at (\cC,  3.6) {陈\,华\,军};
\node (cy)  at (\cC,  1.2) {陈\hspace{1em}永\\杜思桂};
\node (cb)  at (\cC, -1.2) {陈\hspace{1em}斌\\宋爱平};
\node (ct)  at (\cC, -3.6) {陈\hspace{1em}涛\\刘延娇};

% 由婚姻线汇流分枝
\coordinate (junC) at (\cC-1.4, -1.2);
\draw[ln] (union) -- (junC);
\foreach \n in {chj,cy,cb,ct} {
  \draw[ln] (junC) to[out=0,in=180] (\n.west);
}

% 华军两段婚姻：前妻 / 张娟
\node[small] (qq) at (\cC+1.7,  4.3) {（前妻）};
\node[small] (zj) at (\cC+1.7,  2.9) {张\hspace{0.6em}娟};
\draw[ln] (chj.east) to[out=0,in=180] (qq.west);
\draw[ln] (chj.east) to[out=0,in=180] (zj.west);

% ── 第四代：孙辈 ─────────────────────────────────────
\node (xk) at (\cD,  4.7) {陈旭鲲\\刘一佳};
\node (jx) at (\cD,  2.9) {陈佳欣};
\node (xp) at (\cD,  1.2) {陈旭鹏\\孟舒晨};
\node (hy) at (\cD, -0.5) {陈泓妤};
\node (yz) at (\cD, -1.9) {陈映竹};
\node (yf) at (\cD, -2.9) {陈奕霏};
\node (xh) at (\cD, -4.3) {陈旭宏};

\draw[ln] (qq.east) to[out=0,in=180] (xk.west);
\draw[ln] (zj.east) to[out=0,in=180] (jx.west);
\draw[ln] (cy.east) to[out=0,in=180] (xp.west);
\draw[ln] (cb.east) to[out=0,in=180] (hy.west);
\draw[ln] (cb.east) to[out=0,in=180] (yz.west);
\draw[ln] (ct.east) to[out=0,in=180] (yf.west);
\draw[ln] (ct.east) to[out=0,in=180] (xh.west);

% ── 第五代：重孙辈 ──────────────────────────────────
\node (qz) at (\cE, 4.7) {刘沁孜};
\draw[ln] (xk.east) to[out=0,in=180] (qz.west);

% ── 代际题注（横排列于上方） ────────────────────────
\node[draw=none, font=\trueka\footnotesize\itshape, text=inkgray]
  at (\cA, 5.7) {祖\,辈};
\node[draw=none, font=\trueka\footnotesize\itshape, text=inkgray]
  at (\cB, 5.7) {兄\,妹};
\node[draw=none, font=\trueka\footnotesize\itshape, text=inkgray]
  at (\cC, 5.7) {子\,女};
\node[draw=none, font=\trueka\footnotesize\itshape, text=inkgray]
  at (\cD, 5.7) {孙\,辈};
\node[draw=none, font=\trueka\footnotesize\itshape, text=inkgray]
  at (\cE, 5.7) {重\,孙};

\end{tikzpicture}
\end{center}
\end{landscape}



% --------- 版权页 / Colophon ---------
\backmatter
\cleardoublepage
\thispagestyle{empty}
\vspace*{\fill}
\begin{center}
  {\booksansthin\Large 陈炳林回忆录}\\[1em]
  {\bookkai\small 陈炳林\ 口述\quad 陈旭鹏\ 整理}\\[2em]
  {\color{inkfaint}\rule{4em}{0.4pt}}\\[2em]
  {\booksansthin\footnotesize\color{inkgray}
    内部刊印\quad 非卖品\\[4pt]
    二〇二六年\\[2em]
    谨呈家人传阅留念
  }
\end{center}
\vspace*{\fill}

\end{document}
